%%%%%%%%%%%%%%%%%%%%%%%%%%%%%%%%%%%%%%%%%%%
\section{Introduction}
\label{s:intro}
%%%%%%%%%%%%%%%%%%%%%%%%%%%%%%%%%%%%%%%%%%%

Rigid-analytic spaces over non-archimedean fields of mixed or positive characteristic have intricate local and global topology. For example, the \'etale fundamental group of the affinoid unit disc is not finitely generated\footnote{We call a profinite group finitely generated if it is topologically finitely generated, i.e.\ if it admits a finitely generated dense subgroup.}. For a smooth and proper rigid-analytic variety $X$ over $\mathbb{C}_p$, it is known thanks to Scholze \cite{Scholze} that the \'etale cohomology groups $\rH^*(X,\bZ/n\bZ)$ are finite, and it is expected\footnote{We learned of this question from Bogdan Zavyalov.} that its \'etale fundamental group is finitely generated, maybe even finitely presented. This is known to hold if $X$ is the analytification of an algebraic variety, but the general case of this question seems to require genuinely new ideas. 

The `culprit' responsible for the lack of finite generation in the affinoid case, and for the difficulty in the smooth and proper case, is the phenomenon of wild ramification. 
Indeed, a finite \'etale cover of a proper rigid-analytic variety may be wildly ramified along the special fibre of a formal model.
In this paper, we avoid this issue and develop a theory of tame fundamental groups $\pitame(X/K)$ of rigid spaces over a non-archimedean field $K$. We show that this invariant is reasonably well-behaved, most notably that it is finitely generated if $X$ is qcqs and $K$ is algebraically closed or a local field.

This will go in parallel to the story for schemes, which we will now recall briefly as a point of reference. Let $X$ be a connected scheme of finite type over an algebraically closed field $k$. A finite \'etale morphism $f\colon Y\to X$ is tame (or ``tamely ramified at infinity'') if for every $y\in Y$, the extension $k(y)/k(f(y))$ is tamely ramified with respect to every valuation of $k(y)$ which is trivial on $k$ (see \cref{sec:tame pi1} for a review of this and other notions of tameness). Tame covers form a Galois category corresponding to the tame \'etale fundamental group $\pitame(X/k)$, a quotient of the \'etale fundamental group $\pi_1(X)$. Finite generation of $\pitame(X/k)$ has been established in various parts by Grothendieck (if $X$ is a smooth curve \cite[XIII, Corollaire 2.12]{SGA1} or a proper variety \cite[X, Th\'eor\`eme 2.9]{SGA1}), Esnault--Kindler \cite{EsnaultKindler2016:LefschetzTheoremsTamely} (if $X$ admits a projective snc compactification), and in \cref{thm:naked pi1tame fg} of this paper in the general case. Recently, finite presentation of $\pitame(X/k)$ has been proved in \cite{EsnaultShustermanSrinivas2022:FinitePresentationTame} (if $X$ admits a projective snc compactification) and  \cite{LaraSrinivasStix2024:FundamentalGroupsProper} (if $X$ is proper), but remains an open question in general.

\addtocontents{toc}{\SkipTocEntry}
\subsection*{Main results}

Let $K$ be a non-archimedean field, i.e.\ a field complete with respect to a rank one valuation. By a rigid space over $K$ we shall mean an adic space locally of finite type over $\Spa(K)$. The difficulty in giving the right definition of tameness in the context of rigid spaces stems from subtleties surrounding the idea of ``boundary at infinity''. In fact, there are two natural notions of tameness, absolute and relative tameness, and it is the relative one that one should use in order to have hopes for finiteness, as we explain below.

Absolute tameness makes use of the fact that every point of an adic space is equipped with a valuation subring $k(x)^+\subseteq k(x)$ of its residue field. 
By definition (see \cite[Definition~3.3]{Hubner2021:AdicTameSite}), a finite \'etale morphism $f\colon Y\to X$ of adic spaces is tame if for every $y\in Y$, the extension of valued fields $(k(y), k(y)^+)/(k(f(y)), k(f(y))^+)$ is tamely ramified. Using this absolute notion of tame covers, we can define the \emph{absolute} tame fundamental group $\pitame(X)$ of a connected adic space $X$. 

However, the tame fundamental group $\pitame(X)$ of a connected rigid space $X$ over $K$ is not topologically finitely generated in general. For example, if $X$ is the affinoid unit disc over $K = \bC_p$, the surjective specialisation map of étale fundamental groups
\[ 
    \pi_1(X) \la \pi_1(\bA^1), \qquad \bA^1 = \Spec(\overline{\mathbb{F}}_p[T])
\]
factors through $\pitame(X)$, and as $\pi_1(\bA^1)$ is not finitely generated, neither is $\pitame(X)$. This is because, by definition of $X = \Spa(K\langle T\rangle)$, we are only testing tameness at valuations on $K\langle T\rangle$ for which $|T(x)|\leq 1$, or which have a centre on $\bA^1$. At the same time, the wild ramification expressed by the non-finite generation of $\pi_1(\bA^1)$ happens ``at infinity'' of $\bA^1$ and is witnessed by a valuation of $\overline{\mathbb{F}}_p[T]$ with $|T(x)|>1$. We refer to \cref{ex:unit-disc-Cp} for more details. In order to address this disparity, we need to add certain boundary points to $X$ corresponding to the \emph{infinity in the special fibre}. This example motivates our \emph{relative} notion of tameness:

\begin{defiABC}[{see \cref{def:tame-rig}}]
    A finite \'etale morphism of rigid spaces $f\colon Y\to X$ is {\bf tame relative to $K$} if for every $y\in Y$, the field extension $k(y)/k(f(y))$ is tamely ramified with respect to every valuation on $k(y)$ whose valuation ring $V$ satisfies $K^\circ\subseteq V\subseteq k(y)^+$.
\end{defiABC}

If $X$ is quasi-compact and separated, one can rephrase this condition using Huber's universal compactification \cite[\S 5.1]{HuberBook}: $Y\to X$ is tame relative to $K$ if the induced finite \'etale morphism $\ov{Y}\to\ov{X}$ of universal compactifications (relative to $K$)  is tame in the absolute sense defined previously (see \cref{prop:tame-S-compactification}). If $X$ is connected and non-empty, then covers $Y\to X$ which are tame relative to $K$ form a Galois category, leading to the \textit{tame relative to $K$} fundamental group $\pitame(X/K)$. 

This turns out to be a well-behaved invariant. For example, $\pitame(X/K)=1$ if $X$ is the unit disc and $K=\overline{K}$ (see \S\ref{ss:tamepi1-rig-fgfp} for further examples).  The following two theorems are our main results regarding $\pitame(X/K)$.

\begin{thmABC}[{see \cref{thm:fg-pi1-rig}}] \label{thmABC:analytic pi1tame fg}
    Suppose that the tame Galois group of $K$ is finitely generated (e.g.\ $K$ is algebraically closed or it is discretely valued and its residue field is either algebraically closed or finite). Let $X$ be a non-empty connected, quasi-compact and quasi-separated rigid-analytic space over $K$. 
    Then $\pitame(X/K)$ is finitely generated.
\end{thmABC}

\begin{thmABC}[{see \cref{thm:fp-sometimes}}] \label{thmABC:analytic pi1tame fp}
    Suppose that $K$ is algebraically closed and let $X$ be a non-empty connected rigid-analytic space over $K$. Suppose moreover that $X$ admits a strictly semistable formal model $\fX$ such that the smooth locus of the special fibre $\fX_k$ admits a projective snc compactification. Then $\pitame(X/K)$ is finitely presented. 
\end{thmABC}

Finite generation of the tame fundamental group is not only a desirable outcome by itself.
It is also crucial in establishing further properties of the tame fundamental group.
In upcoming work, we aim to prove a Künneth formula and invariance under algebraically closed base field extensions for the tame fundamental group.
A major input in our approach to the Künneth formula is the present result on finite generation.


\addtocontents{toc}{\SkipTocEntry}
\subsection*{Proof strategy}

Before we discuss the contents of the paper in detail, let us briefly explain the overall strategy. The proof of \cref{thmABC:analytic pi1tame fg} starts with a reduction, using desingularisation results due to Temkin and Gabber, to the case $K$ algebraically closed and $X$ admitting a semistable formal model $\fX$. The next step is to compare the tame fundamental group of $X/K$ with a suitably defined tame Kummer \'etale fundamental group $\pitame(\fX^{\log}_k/k)$ of the special fibre $\fX_k$ endowed with its natural log structure. This is difficult to do in part because the log structures in question do not satisfy the customary finiteness conditions, and also because it is unclear why the two notions of ``ramification at infinity'' for $X/K$ and for $\fX^{\log}_k/k$ should be compatible. Having done this step, we are left with showing that $\pitame(Y/k)$ is finitely generated for a certain class of log schemes $Y$ over $k$ with non-finitely generated log structures. Using a ``formal gluing'' type argument along the log stratification, this question is reduced to the case of locally constant log structures. In turn, this case is easily reduced to the case of trivial log structures: the finite generation of $\pitame(Y/k)$ for a scheme of finite type over $k$. Using alterations, finite generation follows from the case of smooth varieties admitting a projective snc compactification, known\footnote{The result form \cite{EsnaultKindler2016:LefschetzTheoremsTamely} exploits Lefschetz theorems for generic hyperplane sections, and, ironically, finally in the case of curves lifts back to characteristic $0$ and appeals to topology via the Riemann Existence Theorem.} thanks to \cite{EsnaultKindler2016:LefschetzTheoremsTamely}. 


\addtocontents{toc}{\SkipTocEntry}
\subsection*{Tame fundamental groups of schemes}

The structure of our paper traces the above argument backwards. In \cref{sec:pit schemes} we treat tame fundamental groups of schemes. For a map of schemes $X\to S$ with $X$ connected and non-empty, one can define the tame fundamental group $\pitame(X/S)$ (\cref{defi:tamepi1 schemes}) by imposing tameness conditions on finite \'etale covers $Y\to X$ (\cref{defi:tame cover of schemes X over S}). Following the constructions in \cite{Temkin2011:RelativeRiemannZariskiSpacesa} and \cite{Hubner2021:AdicTameSite}, this is the tame fundamental group of the discretely ringed adic space $\Spa(X,S)$.
This definition is non-standard in the classical literature on tame covers but should be considered the most natural one.
In \cref{sec:adic tame discrete case} we compare it with existing notions of tameness and also report on improvements on the comparison theorems in \cite{KerzSchmidt2010:DifferentNotionsTameness} due to new developments in non-archimedean geometry, see \cite{HuebnerTemkin2026:WildLocus}.

Van Kampen/descent type arguments for Galois categories play a big role in our paper, and we review the relevant facts in \cref{sec:descent for Gal cats}, in particular, giving a descent criterion for finite generation and presentation of tame fundamental groups (\cref{cor:DD-fgfp}) that essentially goes back to \cite[IX, \S5]{SGA1}. 
These descent techniques find their first application in \cref{thm:naked pi1tame fg} (resp.~\cref{thm:naked pi1tame fp}) to show that the tame fundamental group of a non-empty connected scheme of finite type over an algebraically closed field is finitely generated (resp.~finitely presented under some extra conditions).
Finite generation and presentation of the tame fundamental group of a sufficiently nice variety was previously known due to \cite{EsnaultShustermanSrinivas2022:FinitePresentationTame} but the above described descent techniques allow us to deduce finite generation in full generality and finite presentation under some extra conditions that are far milder than smoothness.

In later sections the descent results from \cref{sec:descent for Gal cats} will also be crucial in treating the tame fundamental groups of log schemes (\cref{s:fg-pi1-log}) and adic spaces (\cref{s:fg-pi1-rig}).

\addtocontents{toc}{\SkipTocEntry}
\subsection*{Tame fundamental groups of log schemes}

In Section~\ref{s:fg-pi1-log} we turn to log schemes.
For later application in \cref{s:fg-pi1-rig} it is essential that we work with monoids that are not necessarily finitely generated.
More precisely, we will encounter log schemes of finite type over $\Spec(M \to k)$, where~$k$ is the residue field of an algebraically closed non-archimedean field $(K,K^\circ)$ and $M = K^\circ \cap K^\times$.
These log schemes are not fs, and accordingly this section relies on the extension of logarithmic geometry beyond fs log schemes in our companion paper \cite{AHLS-Log} (see the review in \S\ref{ss:log-review}), which in particular defines the (Kummer) \'etale fundamental group $\pi_1(X)$ of a saturated log scheme $X$. In this paper, we only need to consider log schemes of type $\typeVd$, i.e.\ those which locally admit a chart by a monoid which is finitely generated over a divisible valuative submonoid. Under this simplifying assumption, every Kummer \'etale map is locally of finite presentation as a map of schemes. In \cref{def:FKEt-pi1-tame} we introduce the tame fundamental group $\pitame(X/S)$ of a log scheme $X$ 
(with underlying scheme $\underline{X}$)
relative to a map of schemes $\underline{X}\to S$ by imposing tameness conditions on the finite (or integral) maps of schemes underlying Kummer \'etale covers. In case the log structure on $X$ is locally constant, we show in \cref{cor:torus-fibration-tame} that we have an exact sequence 
\[
        \begin{tikzcd}
            \Hom\big(\ov \cM^\gp_{X,\bar x} ,\hZ'(1)\big) \ar[r] & \pitame(X/S,\bar x) \ar[r] & \pitame(\underline{X}/S,\bar x) \ar[r] & 1.
        \end{tikzcd}
    \]
Thanks to our $\typeVd$ assumption, the group on the left is a finitely generated $\widehat{\bZ}'$-module (the divisible part of the monoid does not contribute), and we conclude that $\pitame(X/S)$ is finitely generated or finitely presented if and only if $\pitame(\underline{X}/S)$ has this property. If $X$ is a log scheme of type $\typeVd$ over an algebraically closed field $k$, then $X$ admits a dense open $U$ on which the log structure is locally constant. In \S\ref{ss:strat}, we employ this observation and a delicate formal gluing argument to show the following result. 

\begin{thmABC}[{see \cref{thm:pi1tame log is fg}}] \label{thmABC:log pi1tame fg}
    Let $X$ be a non-empty connected log scheme of type $\typeVd$ whose underlying scheme is of finite type over an algebraically closed field $k$.  Then $\pitame(X/k,\bar x)$ is finitely generated.
\end{thmABC}

\noindent 
We also show how $\pitame(X/k)$ can be computed in practice in some situations (smooth over a log point) by performing ``surgery'' that disassembles a log scheme into strata closures with their compactifying log structures (\S\ref{ss:surgery}). This gives sufficient conditions for $\pitame(X/k)$ to be finitely presented (\cref{cor:surgery-fp}), a crucial ingredient in the proof of \cref{thmABC:analytic pi1tame fp}. 

\addtocontents{toc}{\SkipTocEntry}
\subsection*{Vertical compactifications of adic spaces}

Section~\ref{s:compactifications} deals with vertical compactifications of adic spaces and can be read independently of the previous material. The goal is to build, for a map of adic spaces $X\to S$, a space $\Spa(X,S)$ of ``valuations on $X$ with a centre on $S$'' which exists under some mild conditions. If $X\to S$ is a map of analytic adic spaces (for example, $X$ is a rigid space over $K$ and $S=\Spa(K)$), then $\Spa(X,S)$ is the universal compactification constructed by Huber \cite[\S 5]{HuberBook}. If $X_0\to S_0$ is a map of schemes, then $\Spa(X_0^{\rm ad}, S_0^{\rm ad})$ is the discretely ringed adic space $\Spa(X_0,S_0)$ mentioned above. Our key case of interest, the space $\Spa(\fX^{\rm ad}, \Spa(K^\circ))$ for an adic formal scheme $\fX$ over $K^\circ$, interpolates between the two constructions. The space $\Spa(X,S)$ has been constructed by Solodov \cite{Solodov-vertical-compactification} following Huber's approach. We prove additional facts about it, and characterise its points as ``root triples'' $(x^\circ, k(x)^+, x^\succ)$ where $x^\circ\in X$ is a point without vertical generalisations, $k(x)^+$ is a valuation subring of $k(x^\circ)$ contained in its valuation ring $k(x^\circ)^+=k(x^\circ)^\circ$, and where $x^\succ$ is a centre of $k(x)^+$ on $S$ (see \cref{defi:root-triple}). These triples are the test objects used in our definition of tameness (\cref{def:tame-rig}).
We want to emphasize that although the definition of tameness only involves the underlying set of $\Spa(X,S)$, the arguments used in \cref{ss:tameness-comparison} rely on the adic space structure of $\Spa(X,S)$, which is the main novelty of this section compared to \cite{Solodov-vertical-compactification}.

\addtocontents{toc}{\SkipTocEntry}
\subsection*{Tame fundamental groups of rigid spaces} 

In Section~\ref{s:fg-pi1-rig} we finally turn our attention to rigid spaces over a non-archimedean field $K$. We introduce the tame fundamental group $\pitame(X/K)$ by imposing a pointwise tameness condition on root triples (\cref{def:tamepi1-rigid}). We show a variant of the fundamental exact sequence (\cref{prop:weak-fundamental-exact-seq}), whose existence allows us to reduce the proof of \cref{thmABC:analytic pi1tame fg} to the case when $K$ is algebraically closed. 
In \S\ref{ss:alterations-v-descent} we use the fact that surjections between qcqs rigid spaces are of descent for finite \'etale covers combined with  uniformisation theorems due to Gabber \cite{TravauxGabberIX} and Temkin \cite{Temkin2017:AlteredLocalUniformization} to reduce the proof further to the case when $X$ admits a strictly semistable formal model $\fX$. At this point, we are enabled to use log geometry, since $\fX$ endowed with the standard log structure (\cref{prop:log-sm-vertical}) becomes $\fX^{\log}$ which is of type $\typeVd$ and log smooth over $\Spf(K^\circ)^{\log}$.

The two most difficult steps in this
  section are the following. First, we show using approximation techniques and a lemma of Zavyalov (\cref{lem:Zavyalov}) that a finite \'etale map $Y\to X$ which is tame (in the absolute sense) extends uniquely to a Kummer \'etale map $\fY^{\log} \to \fX^{\log}$ (\cref{thm:Abhyankar-formal-more-precise}). In terms of fundamental groups, this means that we have an isomorphism $\pitame(X) \simeq \pi_1(\fX^{\log}_k)$ where $\fX^{\log}_k$ is the log special fibre of $\fX^{\log}$ (see \cref{cor:generic fibre and eta normalisation} and \cref{prop:FEt-formal-vs-sp-fibre}). 
  Second, we show that this isomorphism induces an isomorphism (\cref{prop:tameness-comparison} and \cref{cor:Abhy-infty-equiv})
\[
    \pitame(X/K) \simeq \pitame(\fX^{\log}_k/k)
\]
of the tame quotients.
The argument for this is delicate, relying on the properties of the vertical compactification in \cref{s:compactifications} and a ``cdh uniformisation property'' (\cref{def:cdh-unif}) of the underlying schemes of certain log schemes. Given the above isomorphism and the aforementioned reductions, \cref{thmABC:analytic pi1tame fg} now follows from \cref{thmABC:log pi1tame fg}, and \cref{thmABC:analytic pi1tame fp} from \cref{cor:surgery-fp}. 


%---------------------------------------
\addtocontents{toc}{\SkipTocEntry}
\subsection*{Acknowledgements}
%---------------------------------------

First and foremost, we thank IMPAN at Sopot for excellent working conditions during our stay in September 2023 when this project started. We also thank Heidelberg University for hosting us during an intense few weeks in February 2024. We thank Paul Alexander Helminck, Ben Heuer, Michael Temkin, Alex Youcis, and Bogdan Zavyalov for valuable discussions.


%---------------------------------------
\addtocontents{toc}{\protect\SkipTocEntry}
\subsection*{Funding}
%---------------------------------------

\grants 


%%%---------------------------------------
\addtocontents{toc}{\SkipTocEntry}
\section*{Notation and conventions}
%%%---------------------------------------


\begin{itemize}[leftmargin=*]
    \item 
    If $X$ is a (formal) scheme, or an adic space, we denote by $\FEt_X$ the category of finite \'etale morphisms to $X$. 
    \item 
    Monoids are commutative, in additive notation, and with~$0$.
    \item 
    We denote log schemes by single symbols e.g.~$X$ (not $(X, \cM_X)$) and the underlying scheme is denoted by $\underline{X}$. For log schemes and maps between them, we use adjectives like regular, smooth, \'etale to mean the respective notions in log geometry. In particular, if $X$ is a log scheme of type $\typeVd$, we denote by $\FEt_X$ the category of finite Kummer \'etale maps to $X$, see \cref{def:Kummer-etale-map}. See \S\ref{ss:log-review} for a review of log geometry as used in this paper.

    \item 
    By \emph{locally} we mean working locally in the \'etale topology; in the logarithmic context this means in the strict \'etale topology, i.e.~the \'etale topology after forgetting the log structure. We nevertheless sometimes say \emph{\'etale locally} to stress the use of the \'etale topology.
    
    \item 
    If $X$ is a non-empty connected (log) (formal) scheme or an adic space, we denote by $\pi_1(X,\bar x)$ the fundamental group of the Galois category $\FEt_X$ with respect to the fibre functor induced by a given geometric (log) point $\bar x\to X$. In particular, we never use the notation $\pi_1^{\log}$ for the fundamental group of a log scheme.
    
    \item 
    For a Huber pair $(A,A^+)$ we denote by $A^\circ\subseteq A$ the ring of powerbounded elements and use $\Spa(A)$ as an abbreviation for $\Spa(A,A^\circ)$. 
    
    \item 
    By a \emph{non-archimedean field} we mean a field $K$ which is complete with respect to a rank one valuation. We denote its valuation ring by $K^\circ=K^+$. A \emph{rigid space over $K$} is an adic space locally of finite type over $\Spa(K)$. See \S\ref{ss:adic-spaces} for our conventions regarding adic spaces, rigid spaces, and formal schemes. 
    
    \item 
    For a point $x$ of an adic space $X$, we denote by $k(x)$ the residue field of $\cO_{X, x}$, by $k(x)^+\subseteq k(x)$ the corresponding valuation ring (so that $\cO_{X,x}^+ \subseteq \cO_{X,x}$ is the preimage of $k(x)^+$), and by $k(x)^\succ$ the residue field of $k(x)^+$ (or equivalently of $\cO_{X,x}^+$), which we call the \emph{specialisation field} of $x$. The last convention is not standard, but it seems that the ``residue field of the residue field'' does not have an established name or notation in the literature. 
        
    \item We call a profinite group finitely generated (resp.~presented) if it is topologically finitely generated (resp.~presented).
    
    \item 
    Once we work over a fixed algebraically closed field $k$ (or over a valuation ring $K^\circ$ with residue field $k$), we denote by $\hZ{}'(1) = \varprojlim \mu_n(k)$ the Tate module of $k^\times$. If $p$ is the characteristic exponent of $k$, the group $\hZ{}'(1)$ is isomorphic to $\hZ' = \prod_{\ell\neq p}\bZ_\ell$, the prime-to-$p$ completion of $\bZ$.

    \item 
    The fundamental group of a monoid $P$ is defined  as $\pi_1(P) = \Hom(P^{\rm gp}, \widehat{\bZ}'(1))$ in \cref{defi:pi1-of-monoid}.
    \item 
    For a field $K$ we denote its absolute Galois group by $\pi_1(K)$ where the omitted base point is some chosen algebraic closure $\ov{K}$ of $K$ so that $\pi_1(K) = \Gal(K^\sep/K)$ with the separable closure $K^\sep$ of $K$ in $\ov{K}$. If $(K,K^+)$ is a henselian valued field, then we denote the maximal tamely ramified quotient of the absolute Galois group by $\pitame(K)$.

    \item For a map of schemes $X\to S$, an \emph{snc compactification} is a factorisation $X\to \overline{X}\to S$ where $\overline{X}\to S$ is smooth and proper and where $X\to \overline{X}$ is an open immersion whose complement is the support of an effective Cartier divisor $D\subseteq \ov{X}$ which is snc (simple normal crossings) relative to $S$, i.e.\ Zariski locally on $\overline{X}$, the divisor $D$ is the preimage of the union of coordinate hyperplanes under an \'etale map to $\bA^n_S$.

    \item For an adic formal scheme $\fX$ over $K^\circ$, we denote by $\fX^{\rm log}$ the log formal scheme obtained by endowing $\fX$ with the \emph{standard log structure} $\cM_\fX^{\rm std}=\cO_\fX\times_{\cO_\fX\otimes K} (\cO_\fX\otimes K)^\times$, i.e.\ the one induced by the rigid generic fibre (see \S\ref{ss:log-formal-schemes}). 

\end{itemize}


